\section*{Подробные решения. Вариант 1}

\subsection*{1. Перевод из восьмеричной системы в шестнадцатеричную}
Каждую восьмеричную цифру заменяем тремя битами, затем группируем биты по четыре от точки в обе стороны:
\[
\begin{aligned}
3260.21_8
&=011\,010\,110\,000.010\,001_2\\
&=0110\,1011\,0000.0100\,0100_2\\
&=\boxed{6B0.44_{16}}.
\end{aligned}
\]
Для дробной части справа добавлены два нулевых бита; они не меняют значение. Проверка: $0.21_8=2/8+1/64=17/64=4/16+4/256=0.44_{16}$.

\subsection*{2. Перевод целого десятичного числа}
Делим на 16 с остатком:
\[
2128=16\cdot133+0,\qquad
133=16\cdot8+5,\qquad
8=16\cdot0+8.
\]
Остатки в обратном порядке дают $\boxed{2128_{10}=850_{16}}$. Проверка: $8\cdot16^2+5\cdot16=2128$.

\subsection*{3. Перевод десятичной дроби}
Последовательно умножаем дробную часть на 8:
\[
0.65625\cdot8=5.25\quad\Rightarrow\quad5,\qquad
0.25\cdot8=2\quad\Rightarrow\quad2.
\]
Дробная часть обнулилась, поэтому $\boxed{0.65625_{10}=0.52_8}$. Проверим:
\[
\frac58+\frac{2}{64}=0.65625.
\]

\subsection*{4. Сложение в пятеричной системе}
Складываем справа налево с переносами. В дробных разрядах $4+1=10_5$: пишем 0 и переносим 1; затем $4+1+1=11_5$: пишем 1 и переносим 1 в целую часть. В целой части последовательно получаем $3+2+1=11_5$, $2+3+1=11_5$, $1+3+1=10_5$:
\[
\begin{array}{r@{}l}
  123&.44_5\\
+ 332&.11_5\\ \hline
 1011&.10_5
\end{array}
\]
Ответ $\boxed{1011.10_5}$. Проверка: $123.44_5=38+24/25$, $332.11_5=92+6/25$; сумма $131.2_{10}=1011.10_5$.

\subsection*{5. Сложение 16-битных дополнительных кодов}
Для $A=342_{10}$ получаем $342=156_{16}$, поэтому его 16-битный код --- \texttt{0000 0001 0101 0110} ($0156_{16}$). Для $430_{10}=01AE_{16}$ инвертируем биты положительного числа и прибавляем единицу:
\[
01AE_{16}\ \longrightarrow\ FE51_{16}\ \longrightarrow\ FE52_{16}.
\]
Таким образом, код $B=-430$ равен \texttt{1111 1110 0101 0010}. Складываем 16-битные слова:
\[
0156_{16}+FE52_{16}=FFA8_{16}
=\texttt{1111 1111 1010 1000}.
\]
Для записи кода в восьмеричной системе дополняем \emph{слева} двумя нулями до 18 битов и группируем по три:
\[
\texttt{00 1111 1111 1010 1000}
=\texttt{001 111 111 110 101 000}
=\boxed{177650_8}.
\]
Проверка знакового значения: $342-430=-88$, а $2^{16}-88=65448=FFA8_{16}$; переполнения знакового диапазона нет. Восьмеричная запись выше относится именно к битовому коду, а не к числу $-88$.

\subsection*{6. Вычитание 16-битных чисел с фиксированной точкой}
\enlargethispage{3\baselineskip}
Сначала переведём восьмеричное число: $1011_8=1\cdot8^3+1\cdot8+1=521_{10}$. Значит, $A=-521_{10}$ и $A-B=-521-56=-577_{10}$.

В 16-битном дополнительном коде положительное $521=0209_{16}$; инверсия и прибавление единицы дают код $A$: $FDF7_{16}$. Для вычитания прибавляем дополнительный код $-B=-56$: из $0038_{16}$ получаем $FFC8_{16}$. Тогда
\[
FDF7_{16}+FFC8_{16}=1FDBF_{16}.
\]
Старший перенос за пределы 16 битов отбрасывается. Ответ --- $\boxed{FDBF_{16}}$ (биты \texttt{1111 1101 1011 1111}). Проверка: $2^{16}-577=64959=FDBF_{16}$; $-577$ принадлежит диапазону $[-32768,32767]$.

\subsection*{7. Представление $123.344_{10}$ в IEEE 754 \texttt{float32}}
Знаковый бит равен $S=0$, так как число положительное. Число лежит между $64$ и $128$, поэтому истинный порядок $E=6$. Смещённый порядок $e=E+127=133=10000101_2$.

В нормализованной записи значащая часть содержит 24 бита с ведущей единицей (23 бита сохраняются в дробном поле). Шаг между соседними числами при $E=6$ равен $2^{6-23}=2^{-17}$. Чтобы определить округление, умножим исходное число на $2^{17}$:
\[
123.344\cdot2^{17}
=16166944+\frac{96}{125}.
\]
Поскольку $96/125>1/2$, округляем до целого $16166945$. Его 24-битная запись --- \texttt{111101101011000000100001}: первая единица скрытая, оставшиеся 23 бита образуют поле \texttt{11101101011000000100001}. Случай равного удаления от соседних чисел здесь не возникает.

Поля: знак $S=0$, порядок $e=10000101_2$.
\par\noindent Дробное поле: \texttt{11101101011000000100001}. Группируем слово по четыре бита:
\[
\texttt{0100 0010 1111 0110 1011 0000 0010 0001}.
\]
Получаем $\boxed{42F6B021_{16}}$. Проверка: код представляет число $123.34400177\ldots$, ближайшее к $123.344$ в \texttt{float32}.

\subsection*{8. Таблица истинности схемы}
Верхний элемент ИЛИ получает $A$ и $B$, исключающее ИЛИ --- $A$ и $C$, нижний НЕ --- $B$. После двух элементов И выход инвертируется:
\[
F=\lnot\bigl((A\lor B)\land(A\oplus C)\land\lnot B\bigr).
\]
\begin{center}
\begin{tabular}{ccc|ccc|c}
$A$&$B$&$C$&$A\lor B$&$A\oplus C$&$\lnot B$&$F$\\\hline
0&0&0&0&0&1&1\\
0&0&1&0&1&1&1\\
0&1&0&1&0&0&1\\
0&1&1&1&1&0&1\\
1&0&0&1&1&1&0\\
1&0&1&1&0&1&1\\
1&1&0&1&1&0&1\\
1&1&1&1&0&0&1
\end{tabular}
\end{center}
Только при $(A,B,C)=(1,0,0)$ все входы последнего И равны единице, поэтому выходной НЕ даёт ноль.
